% =====================================================================
% SUPPLEMENTAL MATERIAL — uploaded as ONE file, "Supplemental File for Review".
%
% RSNA wants all supplemental text, tables and figures in a single document, and
% does not copyedit it. Everything here exists because the 3000-word Original
% Research limit will not hold it, not because it is unimportant: a reviewer
% should not have to reverse-engineer the repository to understand the
% experiment.
%
% Every number is regenerated by ../scripts/reproduce.sh from ../results/.
% =====================================================================
\documentclass[11pt]{article}
\usepackage[margin=1in]{geometry}
\usepackage{booktabs}
\usepackage{amsmath}
\usepackage[hidelinks]{hyperref}
\usepackage{setspace}
\usepackage{times}
\doublespacing
\raggedright
\pagestyle{empty}
\setlength{\parindent}{0pt}
\setlength{\parskip}{0.6em}
\renewcommand{\thesection}{S\arabic{section}}
\renewcommand{\thetable}{S\arabic{table}}

\begin{document}

{\large\bfseries Supplemental Material\par}

Supporting ``Supervised Tumor Attention Transfers Across CT Cohorts, but
\textit{EGFR} Prediction Does Not Exceed Clinical Variables.''

\section{Training configuration}

Reported so the experiment can be understood without reading the repository.
All values are the locked configuration used for every run in the manuscript.

\begin{tabular}{ll}
\toprule
Parameter & Value \\
\midrule
Encoder & BiomedCLIP ViT-B/16, last two transformer blocks unfrozen \\
Region head & 8 learned queries, cross-attention to patch tokens \\
Optimizer & AdamW \\
Learning rate & $2\times10^{-4}$, backbone multiplier 1.0 \\
Schedule & OneCycle, warmup ratio 0.05 \\
Weight decay & 0.01 \\
Batch size & 8, gradient accumulation 2 (effective 16) \\
Epochs & 40 maximum \\
Early stopping & patience 8, minimum 10 epochs \\
Gradient clipping & 1.0, with non-finite steps skipped \\
Loss weights & classification 1.0, attention 0.5, generation 0.0 \\
Model selection & composite inner-validation monitor: \\
 & \quad $1.0\times$ \textit{EGFR} AUC $+\ 0.5\times$ attention-mass lift \\
Sections per volume & 16 \\
Crop & $(\text{bbox}_{\text{mm}} + 24)\times2.5$, 30\% positional jitter \\
Window & center $-600$~HU, width 1500~HU \\
Input size & $224\times224$ \\
Hardware & Apple silicon, Metal Performance Shaders backend \\
\bottomrule
\end{tabular}

The generation loss weight is zero in every reported run. With that term active,
folds hit a marginal overflow in the backward pass on this backend, skip every
optimizer step and report an untrained model; the language decoder is therefore
present in the code but contributes nothing, and the model is image-only.
\texttt{gen\_loss} is correspondingly omitted from the selection monitor, because
monitoring an untrained term selects on noise.

No hyperparameter search was performed. These values were fixed before the runs
reported here and were not tuned against any reported endpoint.

\section{Participant flow and selection}

\subsection{Exclusions, with reasons}

\begin{tabular}{lrl}
\toprule
Stage & $n$ & Reason for loss \\
\midrule
Clinical sheet & 211 & --- \\
With an imaging series & 188 & 23 had no CT series in the collection \\
Successfully preprocessed & 158 & 30 had no usable tumor localization: no \\
 & & segmentation and no annotation point, a point \\
 & & unresolvable to a section, or a point landing \\
 & & in air ($<-700$~HU) \\
Known \textit{EGFR} call & 153 & 5 recorded as Unknown or Not collected \\
Adenocarcinoma & 133 & 20 squamous or not otherwise specified \\
With a segmentation & 97 & 36 had an annotated point but no pixel mask \\
\bottomrule
\end{tabular}

\subsection{Analyzed versus excluded patients}

The 58 patients not analyzed are similar to the 153 analyzed on age, sex and
smoking, and differ materially on histology.

\begin{tabular}{lcc}
\toprule
Characteristic & Analyzed ($n=153$) & Excluded ($n=58$) \\
\midrule
Age (y), mean $\pm$ SD & 68.8 $\pm$ 8.8 & 65.9 $\pm$ 12.5 \\
Pack-years, mean $\pm$ SD & 40.5 $\pm$ 26.4 & 37.0 $\pm$ 25.9 \\
Male & 64.1\% & 63.8\% \\
Never-smoker & 21.6\% & 25.9\% \\
Adenocarcinoma & 86.9\% & 67.2\% \\
Squamous cell carcinoma & 11.1\% & 31.0\% \\
\bottomrule
\end{tabular}

\textbf{This is a real selection effect and we do not minimize it.} Exclusion
removed squamous patients at nearly three times the rate of the analyzed cohort.
Because every squamous patient in this collection is \textit{EGFR} wild-type, the
exclusions are enriched for patients whose genotype is inferable from histology
alone. The direction of that bias is to make the analyzed cohort \textit{harder}
than the full collection, so it does not inflate the reported discrimination ---
but a study reporting higher performance on a cohort selected this way would be
partly reporting its selection.

\section{The six analytic choices}

The manuscript states that the radiomics comparator moves 0.136 in AUC across six
defensible analytic choices. Those choices are two feature definitions crossed
with three regularization penalties, on identical patients and folds. The
clinical model is shown under the same three penalties for comparison: it moves
0.010.

\begin{tabular}{llcc}
\toprule
Features & Penalty & AUC & 95\% CI \\
\midrule
Legacy, 18 features & L2 & 0.662 & 0.568, 0.754 \\
Legacy, 18 features & L1 & 0.605 & 0.501, 0.706 \\
Legacy, 18 features & L1 + L2 & 0.604 & 0.498, 0.707 \\
Rich, 61 features & L2 & 0.604 & 0.497, 0.707 \\
Rich, 61 features & L1 & 0.526 & 0.416, 0.631 \\
Rich, 61 features & L1 + L2 & 0.568 & 0.457, 0.671 \\
\midrule
\multicolumn{2}{l}{\textit{Radiomics span}} & \textbf{0.136} & \\
\midrule
Clinical, 5 variables & L2 & 0.774 & 0.677, 0.859 \\
Clinical, 5 variables & L1 & 0.773 & 0.675, 0.861 \\
Clinical, 5 variables & L1 + L2 & 0.764 & 0.665, 0.853 \\
\multicolumn{2}{l}{\textit{Clinical span}} & \textbf{0.010} & \\
\bottomrule
\end{tabular}

A separate comparison holds features and estimator fixed and varies only the
validation protocol: flat versus nested cross-validation moves the same arm by
0.066. Neither of these is detectable from an internal point estimate, which is
the point.

\section{Reference standard}

\textit{EGFR} status is distributed with the collection and was determined by
sequencing of surgical or biopsy tissue by the source investigators. The
collection does not publish the sequencing platform, the assay panel, the
per-patient tissue source, or the variant-level calls, and we did not have access
to them. We therefore report status as the binary label the collection provides
and state here that the assay, panel and specimen type could not be harmonized or
verified per patient. Any study using this collection is subject to the same
limit, and one consequence is that a variant-level analysis --- separating
canonical activating variants from exon-20 insertions or passenger mutations ---
is not possible on these data.

The interval from index test to reference standard has a median of 38 days
(IQR, 19.5--65.5). No accrual window is recoverable: the collection's CT dates
span 1989--1996 and are de-identification offsets, so a range derived from them
would misdescribe the cohort. The interval survives the offset because it is a
difference of two offset dates.

The reference standard for attention is the expert tumor segmentation
distributed with each collection, one contour per lesion. Interreader agreement
therefore cannot be computed. A mask perturbation sweep (dilation, erosion,
translation) is reported as robustness and is not a substitute for a second
reader.

\section{Multiplicity and endpoint thresholds}

\subsection{The two co-primary endpoints}

H\textsubscript{1} (external attention alignment) and H\textsubscript{2}
(incremental AUC over the clinical model) test different constructs on different
populations --- 420 external patients and 153 internal patients respectively ---
and were both declared before external evaluation. Each was evaluated at
$\alpha = .05$ without a gatekeeping rule between them, which is the design
actually run and is stated here rather than reconstructed as a hierarchy after
the fact. Because H\textsubscript{2} was pre-declared as an expected null and the
observed difference is negative, no multiplicity adjustment between the two
endpoints would change any conclusion in the manuscript.

Holm correction is applied within the ten-arm pooling sweep, where it does change
a conclusion: the best arm moves from $P_{\text{raw}} = .03$ to
$P_{\text{Holm}} = .28$, and the peritumoral result is reported as suggestive
rather than established because of it.

\subsection{Origin of the 71.4\% fold-rate threshold}

H\textsubscript{1} was pre-specified as ``beats the per-fold spatial shuffle in at
least 10 of 14 evaluations,'' a count taken from the fold total of the internal
run set at the time the plan was fixed. The number of external evaluations is
however many run $\times$ fold checkpoints exist, which for the canonical three
runs is 15. The count was therefore expressed as the same proportion,
$10/14 = 71.4\%$, to preserve the intent rather than silently gaining or losing a
required fold. This is recorded as a deviation in the analysis plan's deviation
log, and the observed 11 of 15 clears it by one evaluation.

\subsection{What ``pre-specified'' does and does not mean here}

The plan, including both endpoints and all controls, was fixed and
version-controlled before external evaluation. It was not fixed before the
external data existed: the cohort was on disk, and the metric's qualitative
behavior on it partly known, when the threshold was chosen. H\textsubscript{1} is
therefore a locked evaluation following preliminary development, not a blind
confirmatory external validation, and we describe it that way. The distinction
matters and is easy to blur, so it is stated in the manuscript's limitations as
well as here.

\end{document}
